\input{build/foundations-preamble.tex}
\usepackage[authoryear]{natbib}
\renewcommand{\refname}{સંદર્ભગ્રંથો}
\renewcommand{\partname}{ભાગ}
\newenvironment{intro}{}{}
\newcommand{\Id}[1]{\mathord{\mathrm{Id}_{#1}}}
\newcommand{\equivrep}[2]{[#1]_{#2}}
\newcommand{\equivclass}[2]{#1/_{\!{#2}}}
\newcommand{\funrestrictionto}[2]{#1\mathord{\restriction}_{#2}}
\newcommand{\funimage}[2]{#1[#2]}
\newcommand{\emptyseq}{\Lambda}
\newcommand{\liff}{\mathbin{\leftrightarrow}}
\newcommand{\dom}[1]{\mathrm{dom}(#1)}
\newcommand{\ran}[1]{\mathrm{ran}(#1)}
\newcommand{\comp}[2]{#2\circ #1}
\newcommand{\pto}{\mathrel{\ooalign{\hfil$\mapstochar\mkern5mu$\hfil\cr$\to$}}}
\newcommand{\fdefined}{\downarrow}
\newcommand{\fundefined}{\uparrow}
\newcommand{\sourcecorrection}[2]{%
  \par\smallskip\noindent
  \fbox{\begin{minipage}{0.94\linewidth}\small
    \textbf{સ્રોત-સુધારો #1.} #2
  \end{minipage}}%
  \par\smallskip}
\definecolor{oldiagcolorD}{HTML}{1973ba}
\newlength{\olphotowidth}
\setlength{\olphotowidth}{.4\textwidth}
\RenewDocumentCommand{\olasset}{O{\olphotowidth} m}{\centering\resizebox{#1}{!}{\input{upstream/#2}}}
\input{build/functions-available-labels.tex}
\renewcommand{\oliflabeldef}[3]{\ifcsname guavailable@#1\endcsname#2\else#3\fi}
\hypersetup{pdftitle={ગણો, સંબંધો અને વિધેયો — ઓપન લોજિક ગુજરાતી}}
\begin{document}
\begin{titlepage}
{\Huge\bfseries ગણો, સંબંધો અને વિધેયો\par}
\vspace{1em}
{\Large ઓપન લોજિક — ગુજરાતી\par}
\vspace{2em}
ઓપન લોજિક પ્રોજેક્ટના અંગ્રેજી મૂળનો ગુજરાતી અનુવાદ.\par
\vspace{1em}
આ આવૃત્તિમાં ગણો, સંબંધો અને વિધેયોનાં સંપૂર્ણ પ્રકરણો છે.
ગણો અને સંબંધોના પાયા પછી, વિધેયોના પ્રકારો, પ્રતિવિધેયો,
સંયોજન અને આંશિક વિધેયોનો અભ્યાસ થાય છે.\par
\vfill
આ યંત્ર દ્વારા કરેલો અનુવાદ છે. સંપૂર્ણ ગ્રંથનું કામ ચાલુ છે.
આ આવૃત્તિમાં મૂળનાં ૨૩ એકમોનો સમાવેશ છે; કુલ ૭૨૨ એકમો છે.
કેટલીક વિશિષ્ટ પરિભાષા કામચલાઉ છે. ચકાસણીની વિગતો અને
મૂળ પાઠ વિશેની નોંધો સંપાદકીય માહિતીમાં આપેલી છે.\par
\vspace{1em}
\href{https://github.com/OpenLogicProject/OpenLogic}{Open Logic Project}
\quad \href{https://creativecommons.org/licenses/by/4.0/}{CC BY 4.0}\par
\href{https://github.com/KokunoYumeto/OpenLogic-translations}{અનુવાદોનું કેન્દ્ર}\par
\end{titlepage}
\tableofcontents
\clearpage
\input{build/functions-body.tex}
\clearpage
\section*{સંપાદકીય નોંધો}
\addcontentsline{toc}{section}{સંપાદકીય નોંધો}
આ નોંધો મૂળ પાઠથી અલગ છે. મૂળની ગણિતીય નિશાનીઓ જાળવી છે.
સંબંધોને ગણ તરીકે રજૂ કરતા ઉદાહરણમાં વિકર્ણ માટે \(I\) વપરાયું છે,
પણ તેને અલગથી વ્યાખ્યાયિત કરેલું નથી. વૃક્ષની શાખાની વ્યાખ્યામાં
\(X\) વપરાયું છે, જ્યારે વૃક્ષનો આધારગણ \(A\) છે; અહીં મૂળનો
\(X\) જાળવ્યો છે. ઉપવૃક્ષના ઉદાહરણમાં ખાલી ઉપગણનો અપવાદ
સ્પષ્ટ નથી, જોકે મૂળવાળા વૃક્ષની વ્યાખ્યા તેને બાકાત રાખે છે.
ક્રમોના વિભાગમાં \(R^+\) સ્વવાચક સંવરણ દર્શાવે છે અને
સંબંધોની ક્રિયાઓના વિભાગમાં તે પરંપરિત સંવરણ દર્શાવે છે;
દરેક વિભાગની સ્થાનિક વ્યાખ્યા લાગુ પડે છે.

વિધેયોના પાયાના વિભાગમાં મૂળનું એક ઉદાહરણ આગત માટે \(n\)
કહ્યા પછી \(x\) વાપરે છે; અનુવાદમાં \(x\) સુસંગત રીતે વાપર્યું છે.
વર્ગમૂળના ઉદાહરણમાં પ્રદેશમાં શૂન્ય હોવાથી “ધન”ને બદલે
“અનૃણ (મુખ્ય)” લખ્યું છે. વિધેયોના આલેખ વિશેના વર્ણનમાં
“\(A\times B\) પરનો સંબંધ”ને બદલે \(A\times B\)માં સમાયેલો
સંબંધ લખ્યો છે. દરેક ફેરફાર પાસે OLFUN ઓળખવાળી અલગ નોંધ છે.

પ્રતિવિધેયોના વિભાગના પ્રથમ વિધાનમાં પ્રદેશ અખાલી હોવાની
છૂટી ગયેલી શરત ઉમેરી છે. ખાલી ગણથી એક ઘટકવાળા ગણમાંનું
વિધેય એક-એક છે, પરંતુ તેનો ડાબી બાજુનો વ્યસ્ત નથી; સાબિતીમાં
\(a\in A\) પસંદ કરવાનો તબક્કો પણ અખાલી પ્રદેશ ધારે છે.
સંબંધ-મર્યાદન અને વિધેય-મર્યાદન વચ્ચેનો બંને અવયવો સામે
માત્ર પ્રદેશનો તફાવત પણ અનુવાદમાં સ્પષ્ટ કર્યો છે.

આવૃત્તિમાં સમાવેલા વિભાગો માટેના સશરતી ઉલ્લેખો સક્રિય છે.
પછીનાં પ્રકરણોના ઉલ્લેખોમાં મૂળે આપેલો વૈકલ્પિક પાઠ વપરાયો છે.
પરિભાષા, મૂળ સ્રોતો અને ચકાસણીની વધુ વિગતો આવૃત્તિની
સાથેની ફાઇલોમાં છે.
\begin{thebibliography}{1}
\bibitem[Benacerraf(1965)]{Benacerraf1965}
Benacerraf, Paul. 1965.
\emph{What numbers could not be}.
The Philosophical Review 74(1), 47--73.
\end{thebibliography}
\end{document}
